% Sources: docs/design_spec.md, apps/fwdbase/README-UDP.md, apps/blockbench/README.md.
% Allocator lineage: eval/2026-11-21-opaque-loop/LEDGER.md, L8.
% Block pilots: eval/2026-11-22-opaque-block-pilot{2,3}. Assistance: original kernel cover letter.
\section{Implementation}
\label{sec:impl}

The Linux prototypes put handle lookup, retention and release in an \iouring
object store. Each endpoint supplies the code that turns incoming bytes into an
object or submits its pages for output. TCP exercises both access modes. The UDP
and raw block-device prototypes use the named interface at different stages of
development.

\paragraph{Sockets.} The TCP collector claims exclusive access to a socket's
receive stream and obtains its bytes through \code{tcp\_read\_sock}. Sending an
object passes its page ranges to the existing socket splice path. A separate
networking change keeps eligible ranges from large pages together on this path,
which also serves \splice and \sendfile. For mapped receive, we modified the
virtio-net driver to draw buffers from the dedicated area and snapshot device
control fields. The driver and stack together keep the headers used for parsing
in private memory (\S\ref{sec:mapped:segregation}). Tests on virtio-net verified
payload delivery through both modes and exercised release of mapped receive
references.

The IPv4 UDP relay has a copying backend and a named backend that targets a
development kernel interface for receiving and sending whole datagrams. The
copying path has run on loopback. The named relay's tests simulate kernel
completions. Its native end-to-end check and the integration of UDP capture with
the mapped receive area remain open.

\paragraph{Block devices.} A separate prototype implements object reads and writes
within a bounded region of a raw device. Device pilots exercised both serving data
over TCP and writing received data back to the device, with payload verification
on reads and readback on writes. These tests used named objects. Mapped block I/O
requires buffers from a dedicated area that the device can use for direct I/O
(\S\ref{sec:named:block}).

\paragraph{Recycling on each CPU.} Object entries often become free on a different
CPU from the one that took them. The store keeps small batches of free entries,
or \emph{magazines}~\cite{magazines}, on each CPU. A task briefly disables
preemption to use its pair of magazines. A software interrupt on the same CPU
uses a separate pair, so local allocations and returns need neither a lock nor
interrupt masking. Hardware interrupt handlers, including non-maskable interrupts,
cannot allocate entries and defer returns.
The shared pool takes a lock when exchanging batches or supplying the last few
entries. The network receive-page pool uses one pair per CPU and briefly disables
software interrupts while accessing it.

Early fixed-pool builds kept free entries per registration. A later kernel build
includes the separate task and interrupt caches. The earlier TCP comparisons in
\S\ref{sec:eval} therefore cannot quantify this allocator's effect. Preallocation
covers object metadata. Sockets, block requests and drivers can still allocate
memory.

\paragraph{Development tools.} OpenAI Codex, GPT-6.1 through OpenCode, GLM and
Anthropic Claude assisted with code, tests, patch descriptions and review for the
named kernel series.
\todod{Confirm the scope of the authors' disclosure for the other components and
manuscript.}
